\documentclass[12pt, a4paper]{article}

% ---------------------------------------------------------------
% Paquetes
% ---------------------------------------------------------------
\usepackage[spanish]{babel}
\usepackage[utf8]{inputenc}
\usepackage[T1]{fontenc}
\usepackage[margin=2.5cm]{geometry}
\usepackage{hyperref}
\usepackage{xcolor}
\usepackage{listings}
\usepackage{enumitem}
\usepackage{fancyhdr}
\usepackage{titlesec}
\usepackage[skins, breakable]{tcolorbox}
\usepackage{booktabs}
\usepackage{graphicx}
\usepackage{parskip}

% ---------------------------------------------------------------
% Colores institucionales
% ---------------------------------------------------------------
\definecolor{uccazul}{RGB}{0, 70, 127}
\definecolor{uccgris}{RGB}{80, 80, 80}
\definecolor{alertanaranja}{RGB}{230, 126, 34}
\definecolor{okverde}{RGB}{39, 174, 96}
\definecolor{codegris}{RGB}{245, 245, 245}

% ---------------------------------------------------------------
% Estilo de código
% ---------------------------------------------------------------
\lstset{
  backgroundcolor=\color{codegris},
  basicstyle=\ttfamily\small,
  breaklines=true,
  frame=single,
  framesep=5pt,
  rulecolor=\color{gray!40},
  xleftmargin=10pt,
  xrightmargin=10pt,
}

% ---------------------------------------------------------------
% Cajas de aviso
% ---------------------------------------------------------------
\newtcolorbox{importante}{
  colback=alertanaranja!10, colframe=alertanaranja,
  title={\bfseries Importante}, fonttitle=\bfseries,
  breakable
}
\newtcolorbox{nota}{
  colback=uccazul!8, colframe=uccazul,
  title={\bfseries Nota}, fonttitle=\bfseries,
  breakable
}
\newtcolorbox{exito}{
  colback=okverde!10, colframe=okverde,
  title={\bfseries Correcto}, fonttitle=\bfseries,
  breakable
}

% ---------------------------------------------------------------
% Encabezado y pie de página
% ---------------------------------------------------------------
\pagestyle{fancy}
\fancyhf{}
\fancyhead[L]{\color{uccgris}\small Sistema de Control de Acceso UCC}
\fancyhead[R]{\color{uccgris}\small Manual de Instalación}
\fancyfoot[C]{\color{uccgris}\small \thepage}
\renewcommand{\headrulewidth}{0.4pt}

% ---------------------------------------------------------------
% Títulos de sección
% ---------------------------------------------------------------
\titleformat{\section}
  {\Large\bfseries\color{uccazul}}{\thesection.}{0.5em}{}[\titlerule]
\titleformat{\subsection}
  {\large\bfseries\color{uccazul!80}}{\thesubsection.}{0.5em}{}

% ---------------------------------------------------------------
% Hipervínculos
% ---------------------------------------------------------------
\hypersetup{
  colorlinks=true,
  linkcolor=uccazul,
  urlcolor=uccazul,
  pdftitle={Manual de Instalacion - UCC Control de Acceso},
  pdfauthor={Corporacion Universitaria Comfacauca}
}

% ---------------------------------------------------------------
% Documento
% ---------------------------------------------------------------
\begin{document}

% === Portada ===
\begin{titlepage}
  \centering
  \vspace*{3cm}
  {\color{uccazul}\rule{\textwidth}{2pt}}\\[1em]
  {\Huge\bfseries\color{uccazul} Sistema de Control de Acceso}\\[0.5em]
  {\LARGE\color{uccgris} Corporación Universitaria Comfacauca}\\[0.5em]
  {\color{uccazul}\rule{\textwidth}{2pt}}\\[2em]

  {\Large\bfseries Manual de Instalación en Equipo Nuevo}\\[0.5em]
  {\large\color{uccgris} Versión 1.0 — 2026}

  \vfill
  {\color{uccgris}
    \textbf{Para uso del personal técnico o administrativo.}\\
    No se requieren conocimientos de programación.
  }
  \vspace{2cm}
\end{titlepage}

% === Tabla de contenidos ===
\tableofcontents
\newpage

% ===============================================================
\section{Introducción}
% ===============================================================

El Sistema de Control de Acceso UCC permite registrar y verificar
el ingreso de estudiantes, empleados y contratistas a las instalaciones
de la universidad. El sistema funciona en una red WiFi local y es
accesible desde cualquier computador o celular conectado a esa red.

\subsection{Arquitectura del sistema}

El sistema consta de cuatro componentes que deben estar ejecutándose
simultáneamente:

\begin{center}
\begin{tabular}{@{}lll@{}}
  \toprule
  \textbf{Componente} & \textbf{Puerto} & \textbf{Función} \\
  \midrule
  PostgreSQL  & 5432 & Base de datos (instala como servicio de Windows) \\
  PostgREST   & 3001 & API REST sobre PostgreSQL \\
  Proxy       & 3000 & Intermediario con manejo de CORS \\
  Frontend    & 80   & Interfaz web accesible por navegador \\
  \bottomrule
\end{tabular}
\end{center}

PostgreSQL se inicia automáticamente con Windows. Los otros tres
componentes los arranca el archivo \texttt{INICIAR.bat}.

% ===============================================================
\section{Requisitos previos}
% ===============================================================

Antes de instalar el sistema, el equipo debe tener instalados
los siguientes programas. Ambos son gratuitos.

\subsection{PostgreSQL 18}

\begin{enumerate}
  \item Descarga el instalador desde:\\
        \url{https://www.postgresql.org/download/windows/}
  \item Selecciona la versión \textbf{18.x} para Windows x86-64.
  \item Ejecuta el instalador. Cuando pida contraseña para el usuario
        \texttt{postgres}, anótala; la necesitarás en el paso 5 de la
        instalación del sistema.
  \item Deja marcado el componente \textbf{PostgreSQL Server}.
        El puerto por defecto (5432) no debe cambiarse.
  \item Finaliza la instalación. Al terminar, selecciona
        \emph{"No lanzar Stack Builder"}.
\end{enumerate}

\begin{importante}
  Recuerda la contraseña que defines para el usuario \texttt{postgres}.
  Si la olvidas, deberás reinstalar PostgreSQL.
\end{importante}

\subsection{Node.js LTS}

\begin{enumerate}
  \item Descarga el instalador desde:\\
        \url{https://nodejs.org/}
  \item Selecciona la versión \textbf{LTS} (la recomendada).
  \item Ejecuta el instalador con las opciones por defecto.
  \item En la pantalla de \emph{"Tools for Native Modules"}, NO es
        necesario marcar la casilla adicional.
\end{enumerate}

\subsection{PostgREST (archivo .exe)}

\begin{enumerate}
  \item Descarga el archivo ZIP desde:\\
        \url{https://github.com/PostgREST/postgrest/releases/tag/v12.2.0}
  \item Busca el archivo: \texttt{postgrest-v12.2.0-windows-x86-64.zip}
  \item Ábrelo y extrae el archivo \texttt{postgrest.exe}.
  \item \textbf{No lo copies aún} — el asistente de configuración te indicará
        la carpeta exacta donde colocarlo en el paso 3 de la instalación.
\end{enumerate}

% ===============================================================
\section{Instalación del sistema}
% ===============================================================

\subsection{Obtener los archivos del sistema}

Los archivos del sistema se entregan como una carpeta comprimida
(\texttt{ucc-control-acceso.zip}) o directamente como carpeta.
Copia esta carpeta a una ubicación estable en el equipo, por ejemplo:

\begin{lstlisting}
C:\ucc-control-acceso\
\end{lstlisting}

\begin{importante}
  No pongas la carpeta en el Escritorio si el usuario del equipo cambia.
  Prefiere \texttt{C:\textbackslash{}ucc-control-acceso\textbackslash{}} o
  \texttt{C:\textbackslash{}Sistemas\textbackslash{}ucc-control-acceso\textbackslash{}}.
\end{importante}

La estructura de la carpeta debe verse así:

\begin{lstlisting}
ucc-control-acceso\
  CONFIGURAR.bat        <- ejecutar primero (solo una vez)
  INICIAR.bat           <- ejecutar cada dia
  postgrest\
    proxy.js
    postgrest.exe       <- lo copias aqui en el paso 3
  database\
    supabase\
      *.sql
  dist\
    index.html
    ...
\end{lstlisting}

\subsection{Ejecutar el asistente de configuración}

\begin{enumerate}
  \item Abre la carpeta del sistema en el Explorador de archivos.
  \item Haz \textbf{clic derecho} sobre el archivo \texttt{CONFIGURAR.bat}.
  \item Selecciona \textbf{"Ejecutar como administrador"}.
  \item Se abrirá una ventana negra. Sigue las instrucciones en pantalla.
\end{enumerate}

\medskip
El asistente realiza los siguientes pasos automáticamente:

\begin{enumerate}[label=\textbf{[\arabic*]}]
  \item Verifica que Node.js y PostgreSQL estén instalados.
  \item Verifica que \texttt{postgrest.exe} esté en la carpeta correcta.
        Si no está, te indica dónde copiarlo y espera.
  \item Instala el servidor de archivos web (\texttt{serve}).
  \item Crea la base de datos y el usuario del sistema en PostgreSQL.
        Te pedirá la contraseña de \texttt{postgres} que definiste al
        instalar PostgreSQL.
  \item Crea todas las tablas y carga el esquema inicial.
  \item Configura PostgREST y copia los archivos necesarios.
  \item Abre los puertos en el firewall de Windows (80, 3000, 3001).
  \item Crea un acceso directo en el Escritorio.
\end{enumerate}

\begin{exito}
  Al finalizar, el asistente muestra la dirección IP del equipo.
  Guarda esa dirección para generar el código QR.
\end{exito}

% ===============================================================
\section{Uso diario — Iniciar el sistema}
% ===============================================================

Cada vez que se encienda el equipo, el sistema debe iniciarse
manualmente. PostgreSQL arranca automáticamente con Windows, pero
los demás componentes deben lanzarse con el \texttt{INICIAR.bat}.

\subsection{Procedimiento}

\begin{enumerate}
  \item Busca el acceso directo \textbf{"Iniciar UCC Control"} en el
        Escritorio (lo crea el asistente de configuración).
  \item Haz \textbf{clic derecho} $\rightarrow$ \textbf{"Ejecutar como administrador"}.
  \item Se abre una ventana negra que muestra las direcciones del sistema.
  \item Deja esa ventana abierta mientras el sistema está en uso.
  \item Para \textbf{detener} el sistema, cierra esa ventana.
\end{enumerate}

\begin{importante}
  No cierres la ventana del sistema mientras haya usuarios conectados.
  Al cerrarla, el sistema queda inaccesible hasta que vuelvas a ejecutar
  \texttt{INICIAR.bat}.
\end{importante}

\subsection{Verificar que el sistema está funcionando}

Abre un navegador en el mismo equipo y entra a:

\begin{lstlisting}
http://localhost
\end{lstlisting}

Debe cargarse la pantalla de inicio del sistema. Si aparece un error,
revisa la sección~\ref{sec:problemas} de este manual.

% ===============================================================
\section{Configurar dirección IP fija}
\label{sec:ipfija}
% ===============================================================

La dirección IP del equipo puede cambiar cada vez que se reconecta
a la red WiFi. Esto haría inválido el código QR. Para evitarlo,
configura una IP estática.

\subsection{Opción A — IP fija desde Windows (recomendada si no hay soporte de red)}

\begin{enumerate}
  \item Presiona \texttt{Win + X} y abre \textbf{Terminal (Administrador)}
        o \textbf{PowerShell (Administrador)}.
  \item Ejecuta estos tres comandos (reemplaza los valores por los de
        tu red si son diferentes):
\end{enumerate}

\begin{lstlisting}
netsh interface ip set address name="Wi-Fi" static 192.168.44.216 255.255.252.0 192.168.44.1
netsh interface ip set dns name="Wi-Fi" static 8.8.8.8
netsh interface ip add dns name="Wi-Fi" 8.8.4.4 index=2
\end{lstlisting}

\begin{nota}
  Los tres valores del primer comando son:
  \begin{itemize}
    \item \textbf{192.168.44.216} — IP que quieres fijar (la que el equipo
          tenía cuando configuraste el sistema).
    \item \textbf{255.255.252.0} — Máscara de subred (no cambiar si la red
          es la del campus).
    \item \textbf{192.168.44.1} — Puerta de enlace (el router del campus).
  \end{itemize}
  Para verificar la IP actual del equipo, ejecuta \texttt{ipconfig} en una
  terminal y busca la línea \emph{Dirección IPv4} del adaptador Wi-Fi.
\end{nota}

\subsection{Opción B — Reserva DHCP en el router (solución institucional)}

Solicita al administrador de red de la universidad que reserve la
dirección IP para la MAC del equipo. Proporciona los siguientes datos:

\begin{lstlisting}
MAC del equipo  : B0-35-9F-9C-7E-4E
IP a reservar   : 192.168.44.216  (o la IP disponible que asignen)
Descripcion     : PC Control Acceso UCC
\end{lstlisting}

Para obtener la MAC de otro equipo, ejecuta en una terminal:
\begin{lstlisting}
getmac /v
\end{lstlisting}

% ===============================================================
\section{Generar el código QR}
% ===============================================================

El código QR permite a los usuarios acceder al portal de registro
desde sus celulares, sin necesidad de escribir la dirección.

\subsection{URL del portal de usuario}

\begin{lstlisting}
http://<IP-DEL-EQUIPO>/login
\end{lstlisting}

Por ejemplo, si la IP es \texttt{192.168.44.216}:
\begin{lstlisting}
http://192.168.44.216/login
\end{lstlisting}

\subsection{Cómo generar el QR}

\begin{enumerate}
  \item Abre cualquier navegador en cualquier equipo.
  \item Ve a \url{https://qr.io} o \url{https://www.qrcode-monkey.com/es/}
  \item Ingresa la URL del portal: \texttt{http://192.168.44.216/login}
  \item Descarga o imprime el QR generado.
\end{enumerate}

\begin{importante}
  El código QR solo funciona desde dispositivos conectados a la red
  WiFi del campus. No funciona desde fuera del campus.
  Si la IP del equipo cambia, el QR queda inválido. Configura IP fija
  (sección~\ref{sec:ipfija}) para que el QR sea permanente.
\end{importante}

% ===============================================================
\section{Panel de administración}
% ===============================================================

\subsection{Acceso}

El panel de administración es exclusivo para administradores del sistema:

\begin{lstlisting}
http://localhost/admin-login
http://<IP-DEL-EQUIPO>/admin-login
\end{lstlisting}

\subsection{Credenciales por defecto}

\begin{center}
\begin{tabular}{@{}ll@{}}
  \toprule
  \textbf{Campo} & \textbf{Valor} \\
  \midrule
  ID institucional & 000000 \\
  Contraseña & 12345678 \\
  \bottomrule
\end{tabular}
\end{center}

\begin{importante}
  Cambia la contraseña por defecto después del primer ingreso.
  Estas credenciales son conocidas por el instalador del sistema.
\end{importante}

\subsection{Funciones disponibles}

Desde el panel de administración se puede:
\begin{itemize}
  \item Ver los registros de acceso en tiempo real.
  \item Gestionar usuarios (estudiantes, empleados, contratistas).
  \item Crear y gestionar administradores adicionales.
  \item Gestionar semestres académicos.
  \item Exportar reportes.
\end{itemize}

% ===============================================================
\section{Solución de problemas}
\label{sec:problemas}
% ===============================================================

\subsection{El navegador muestra "Esta página no está disponible"}

El sistema no está iniciado. Ejecuta \texttt{INICIAR.bat} como
administrador y espera unos segundos antes de recargar el navegador.

\subsection{Aparece un error "502 Bad Gateway"}

PostgREST no está corriendo. Esto puede ocurrir si:
\begin{itemize}
  \item El sistema se inició pero PostgREST tardó en arrancar. Espera
        5 segundos y recarga.
  \item PostgreSQL no está activo. Ve a \textbf{Servicios de Windows}
        (presiona \texttt{Win + R}, escribe \texttt{services.msc}) y
        verifica que el servicio \textbf{postgresql-x64-18} esté en
        estado \emph{En ejecución}.
\end{itemize}

\subsection{El login dice "Credenciales incorrectas" pero la contraseña es correcta}

Verifica que el sistema esté completamente iniciado (los tres procesos:
PostgREST, proxy y frontend). Cierra el \texttt{INICIAR.bat} y
ejecútalo nuevamente.

\subsection{Otros equipos no pueden acceder}

\begin{enumerate}
  \item Verifica que estén en la misma red WiFi del campus.
  \item Confirma que el firewall esté configurado: ejecuta
        \texttt{CONFIGURAR.bat} nuevamente (es seguro ejecutarlo más de
        una vez) para que reconfiguré las reglas.
  \item Asegúrate de estar usando la IP correcta del equipo servidor.
        Abre una terminal en el equipo servidor y ejecuta \texttt{ipconfig}.
\end{enumerate}

\subsection{Cómo reiniciar completamente el sistema}

\begin{enumerate}
  \item Cierra la ventana de \texttt{INICIAR.bat} si está abierta.
  \item Vuelve a ejecutar \texttt{INICIAR.bat} como administrador.
\end{enumerate}

% ===============================================================
\section{Carga masiva de usuarios}
% ===============================================================

Para cargar listas de estudiantes, empleados o contratistas de forma
masiva, el sistema acepta archivos CSV con un formato específico.

Los archivos de ejemplo se encuentran en:
\begin{lstlisting}
database\
  test_estudiantes.csv
  test_empleados.csv
  test_contratistas.csv
  ejemplo_carga_semestral.csv
\end{lstlisting}

Consulta el archivo \texttt{database\textbackslash{}README\_DATABASE.md} para
ver el formato exacto de cada columna.

% ===============================================================
\section{Respaldo de la base de datos}
% ===============================================================

Se recomienda hacer un respaldo periódico de la base de datos para
no perder los registros de acceso.

\subsection{Crear un respaldo}

Abre una terminal (\textbf{no} es necesario ejecutar como administrador)
y ejecuta:

\begin{lstlisting}
"C:\Program Files\PostgreSQL\18\bin\pg_dump" -U postgres -d ucc_control -f respaldo.sql
\end{lstlisting}

Esto crea el archivo \texttt{respaldo.sql} en la carpeta actual.
Cópialo a un USB o carpeta segura.

\subsection{Restaurar un respaldo}

\begin{lstlisting}
"C:\Program Files\PostgreSQL\18\bin\psql" -U postgres -d ucc_control -f respaldo.sql
\end{lstlisting}

% ===============================================================
\section{Contacto y soporte}
% ===============================================================

Este sistema fue desarrollado para la Corporación Universitaria
Comfacauca. Para soporte técnico, contactar al área de sistemas
o al desarrollador original del proyecto.

\vspace{1cm}
\begin{center}
  {\color{uccazul}\rule{0.5\textwidth}{0.5pt}}\\[0.5em]
  {\small\color{uccgris}
    Sistema de Control de Acceso UCC — Manual de Instalación v1.0 \\
    Corporación Universitaria Comfacauca, 2026
  }
\end{center}

\end{document}
